% =============================================================================
\section{Conclusion}\label{sec:conclusion}
% =============================================================================

We introduced a formal perturbation--variance framework for effective context length (ECL) in embedding-generating sequence models. Our definitions---including the Effective Context Profile (ECP), Effective Context Dimension (ECD), directional ECL, interaction ECL, and Context Decay Spectroscopy (CDS)---provide a multi-faceted, statistically rigorous characterization of how much input context a model actually uses. We established key theoretical properties including monotonicity, invariance, exponential decay bounds, minimax estimation rates, and connections to Sobol sensitivity analysis and Blackwell ordering. We developed a suite of estimators with non-asymptotic concentration guarantees, variance reduction techniques, and formal model comparison tests.

The proposed experimental program---spanning ten protocols across eight genomic models, four perturbation types, three locus classes, and biological ground truth---provides a comprehensive blueprint for empirical ECL assessment. Novel contributions including the genomic Needle-in-a-Haystack (gNIAH), ECL-as-training-diagnostic, and cross-model transfer validation extend the framework beyond measurement toward actionable model development insights.

ECL complements nominal context length by quantifying actual context utilization, enabling principled model comparison and mechanistic interpretation. As genomic language models continue to scale toward chromosome- and genome-length inputs, understanding what they effectively see---not just what they nominally accept---becomes essential for both scientific reliability and computational efficiency.
